\chapter{Conclusiones y Trabajo Futuro}
\label{cap:Conclusiones}
El trabajo desarrollado ha permitido alcanzar el objetivo principal planteado al inicio de la memoria: 
construir un prototipo funcional de apoyo al análisis de imágenes biomédicas preclínicas mediante técnicas de visión por computador e inteligencia artificial.
La herramienta obtenida permite segmentar automáticamente el cerebro y las regiones de isquemia cerebral en imágenes de resonancia magnética de ratón, calcular 
medidas volumétricas asociadas e incorporar mecanismos de explicabilidad para analizar el comportamiento de los modelos.

A partir de los resultados obtenidos, puede concluirse que la aproximación propuesta es adecuada como primera solución automatizada dentro de un flujo de trabajo supervisado.
El sistema no sustituye el criterio del investigador experto, pero sí permite reducir parte de la carga manual asociada a la delimitación de regiones de interés y proporciona 
una salida cuantitativa que puede ser revisada, corregida y utilizada como apoyo en estudios biomédicos.

El modelo de segmentación cerebral ha mostrado un comportamiento especialmente sólido. Los resultados alcanzados, con un coeficiente Dice 3D medio de 0,973 y una IoU 3D media 
de 0,947, indican una alta correspondencia entre las máscaras generadas automáticamente y las segmentaciones manuales de referencia. Además, el bajo error volumétrico obtenido 
confirma que el modelo no solo reproduce adecuadamente la forma del cerebro, sino que también conserva de manera precisa su volumen total. Esta estabilidad convierte la segmentación
cerebral en una base fiable para las fases posteriores del pipeline.

La segmentación de isquemia cerebral ha presentado una mayor dificultad, como era esperable por la naturaleza del problema. Las lesiones isquémicas muestran bordes difusos, 
variaciones de intensidad y zonas de transición cuya delimitación depende en parte del criterio del anotador. Aun así, el modelo ha obtenido resultados coherentes, con un Dice 3D
medio de 0,820 y una IoU 3D media de 0,717. Estos valores, aunque inferiores a los obtenidos en la segmentación cerebral, resultan razonables para una tarea con mayor ambigüedad 
anatómica y visual. El modelo permite localizar las regiones patológicas principales y ofrece estimaciones volumétricas útiles como primera aproximación automatizada.

Desde el punto de vista de la cuantificación, el sistema cumple una función relevante, ya que transforma las máscaras generadas en medidas volumétricas interpretables.
Esta salida resulta especialmente importante en el contexto del trabajo, puesto que el interés del usuario final no se limita a visualizar una segmentación, sino también a 
disponer de valores numéricos que permitan comparar casos, estudiar la evolución de una lesión o analizar diferencias entre sujetos.

La incorporación de técnicas de Inteligencia Artificial Explicable ha aportado una capa adicional de análisis. Los mapas generados mediante Grad-CAM, Grad-CAM++ e 
Integrated Gradients han permitido estudiar qué regiones de la imagen influyen en las predicciones de los modelos. En los casos correctamente segmentados, las explicaciones 
obtenidas se concentran en zonas coherentes con las estructuras segmentadas, lo que refuerza la interpretación de que los modelos no se apoyan únicamente en artefactos o regiones 
externas irrelevantes. No obstante, estas técnicas deben entenderse como una herramienta complementaria de inspección, no como una validación clínica independiente.

En conjunto, el trabajo demuestra que el uso de arquitecturas U-Net con aprendizaje por transferencia es una estrategia viable para abordar tareas de segmentación biomédica
en un contexto de datos limitados. La decisión de emplear dos modelos diferenciados, uno para cerebro y otro para isquemia, ha permitido estructurar el problema de forma clara
y aprovechar la segmentación anatómica como paso previo al análisis de la lesión. Esta organización mejora la coherencia del sistema y facilita su interpretación.

A pesar de estos resultados, el trabajo presenta varias limitaciones. La primera y más importante es el tamaño reducido del conjunto de datos. Aunque esta situación es habitual 
en proyectos de imagen biomédica, limita la capacidad de generalización de los modelos y obliga a interpretar los resultados como una validación exploratoria. Una evaluación más
robusta requeriría incorporar un mayor número de casos, distintos protocolos de adquisición y una mayor variabilidad en la morfología e intensidad de las lesiones.

Otra limitación relevante es la dependencia de las anotaciones manuales utilizadas como referencia. En el caso de la isquemia, las fronteras de la lesión no siempre están
claramente definidas, por lo que el ground truth contiene un cierto grado de subjetividad. Esta variabilidad afecta directamente a métricas como Dice e IoU, que penalizan
pequeñas diferencias espaciales aunque estas puedan ser aceptables desde el punto de vista experto. Por ello, los resultados deben interpretarse teniendo en cuenta que la
segmentación manual representa una referencia operativa, no una verdad absoluta.

También debe considerarse que la validación se ha realizado sobre un número limitado de casos de prueba. Aunque estos casos permiten estudiar situaciones relevantes, como
control, isquemia transitoria e isquemia permanente, no son suficientes para afirmar que el sistema generaliza a todos los escenarios posibles. Antes de su uso en estudios 
experimentales más amplios, sería necesario realizar una validación externa con más sujetos y, preferiblemente, con revisión por parte de varios expertos.

\section{Trabajo Futuro}

Como trabajo futuro, una primera línea de mejora consistiría en ampliar el conjunto de datos disponible. La incorporación de nuevos casos permitiría entrenar modelos más 
robustos, reducir el riesgo de sobreajuste y evaluar con mayor precisión el comportamiento del sistema ante lesiones de distinta morfología, tamaño e intensidad. También 
sería recomendable estudiar técnicas de aumento de datos específicas para imagen médica, siempre que las transformaciones aplicadas mantengan la coherencia anatómica de las imágenes.

Una segunda línea de trabajo sería mejorar la evaluación experimental. En futuras versiones deberían utilizarse particiones de entrenamiento, validación y test más amplias, 
así como validación cruzada y análisis separados por tipo de lesión, tiempo de adquisición y calidad de imagen. Además, sería especialmente útil comparar las predicciones del
modelo con anotaciones realizadas por varios expertos, con el fin de cuantificar la variabilidad interobservador y contextualizar mejor el rendimiento del sistema automático.

Desde el punto de vista técnico, podrían explorarse arquitecturas más avanzadas que la U-Net base empleada en este trabajo, como U-Net++, Attention U-Net o modelos con 
mecanismos de atención. Estas arquitecturas podrían mejorar la detección de lesiones pequeñas o de bordes poco definidos. Sin embargo, su uso debería evaluarse con cautela, 
ya que un aumento de complejidad no garantiza necesariamente una mejora en contextos con pocos datos anotados.

Otra posible mejora sería estudiar modelos tridimensionales. El enfoque actual trabaja corte a corte con imágenes 2D, lo que reduce el coste computacional y facilita el uso
de modelos preentrenados, pero no aprovecha completamente la continuidad espacial entre slices consecutivas. Una arquitectura 3D podría mejorar la coherencia volumétrica de 
las predicciones, aunque requeriría más datos, mayor capacidad computacional y una estrategia de entrenamiento más compleja.

También sería conveniente desarrollar una interfaz de usuario orientada al investigador biomédico. Aunque los notebooks entregados permiten reproducir el entrenamiento, la
inferencia y la evaluación, una aplicación con interfaz gráfica facilitaría la carga de volúmenes, la visualización de cortes, la edición de máscaras, el cálculo de volúmenes
y la exportación de resultados. Esta mejora favorecería la integración de la herramienta en un entorno de trabajo real.

Finalmente, la metodología desarrollada podría extenderse a otros órganos, lesiones o modalidades de imagen. Aunque este trabajo se ha centrado en cerebro e isquemia
cerebral en ratón, el pipeline propuesto proporciona una base adaptable a otros problemas de segmentación biomédica, siempre que se disponga de datos anotados y se ajusten
adecuadamente las fases de preprocesamiento, entrenamiento y validación.

En conclusión, este Trabajo de Fin de Grado ha permitido desarrollar una primera herramienta funcional para el análisis automatizado de imágenes biomédicas preclínicas.
Los resultados obtenidos muestran un rendimiento sólido en segmentación cerebral y un comportamiento prometedor en segmentación de isquemia, especialmente teniendo en 
cuenta la dificultad del problema y la limitación de datos disponibles. Aunque el sistema requiere validación adicional y mejoras antes de su aplicación en entornos experimentales
más amplios, constituye una base técnica válida para avanzar hacia herramientas de apoyo que reduzcan la carga manual, mejoren la reproducibilidad y faciliten la cuantificación de 
lesiones en investigación biomédica.


